\section{Metodología Propuesta}\label{sec:methodology}
 
\subsection{Pipeline general}
El pipeline metodológico propuesto se estructura en las siguientes etapas secuenciales: (1) ingesta y diagnóstico de anomalías; (2) corrección de códigos de error de sensor; (3) ingeniería de la variable objetivo con prevención de \emph{data leakage}; (4) partición estratificada; (5) imputación por mediana y escalamiento robusto ajustados únicamente sobre el conjunto de entrenamiento; (6) reducción de dimensionalidad mediante PCA; (7) entrenamiento y optimización de un Random Forest Classifier balanceado; (8) evaluación integral sobre el conjunto de prueba. La Fig. \ref{fig:pipeline} resume esta arquitectura.
 
\begin{figure}[H]
\centering
\begin{verbatim}
%% DIAGRAMA MERMAID (a convertir a figura IEEE)
%% graph TD
%%   A[Ingesta DATASET.csv] --> B[Diagnostico EDA -999]
%%   B --> C[Correccion NaN]
%%   C --> D[Target_Helada + eliminacion leakage]
%%   D --> E[Split 80/20 estratificado]
%%   E --> F[Imputacion mediana + RobustScaler]
%%   F --> G[PCA K=5, 95% varianza]
%%   G --> H[RandomForest + GridSearchCV]
%%   H --> I[Evaluacion: ROC, PR, Kappa]
\end{verbatim}
\caption{Diagrama de flujo del pipeline metodológico propuesto (representación textual Mermaid, pendiente de conversión a figura vectorial IEEE).}
\label{fig:pipeline}
\end{figure}
 
\subsection{Algoritmo de clasificación}
Se seleccionó un ensamble de árboles de decisión (\emph{Random Forest Classifier}) por su afinidad geométrica con particiones ortogonales del espacio, coherente con la naturaleza del subespacio generado por PCA. El modelo se configuró con ponderación balanceada de clases (\texttt{class\_weight='balanced'}) para compensar el desbalance de clases documentado en la Sección III-G, ajustando los pesos de manera inversamente proporcional a la frecuencia de cada clase:
\begin{equation}
w_c = \frac{N}{2 \cdot N_c}
\label{eq:classweight}
\end{equation}
donde $N$ es el número total de instancias y $N_c$ el número de instancias de la clase $c$.
 
\subsection{Función objetivo y optimización de hiperparámetros}
La optimización se realizó mediante \texttt{GridSearchCV} con validación cruzada de 3 particiones (\emph{3-fold}), utilizando como función objetivo el F1-Score Macro:
\begin{equation}
F1_{macro} = \frac{1}{C} \sum_{c=1}^{C} \frac{2 \cdot P_c \cdot R_c}{P_c + R_c}
\label{eq:f1macro}
\end{equation}
donde $P_c$ y $R_c$ son la precisión y el \emph{recall} de la clase $c$, respectivamente. Esta métrica fue elegida en lugar de la exactitud (\emph{accuracy}), ya que, dado el desbalance del 83.76\% frente a 16.24\%, un clasificador trivial que prediga siempre la clase mayoritaria superaría el 90\% de exactitud sin ninguna utilidad práctica.
 
La Tabla \ref{tab:grid} resume el espacio de búsqueda de hiperparámetros explorado.
 
\begin{table}[H]
\centering
\caption{Espacio de búsqueda de hiperparámetros (\texttt{GridSearchCV})}
\label{tab:grid}
\begin{tabular}{lll}
\toprule
\textbf{Hiperparámetro} & \textbf{Valores explorados} & \textbf{Descripción} \\
\midrule
\texttt{n\_estimators} & \{50, 100\} & Número de árboles \\
\texttt{max\_depth} & \{5, 10\} & Profundidad máxima \\
\texttt{min\_samples\_split} & \{2, 5\} & Mínimo de muestras por nodo \\
\bottomrule
\end{tabular}
\end{table}
 
La grilla resultante comprende 8 combinaciones de hiperparámetros, evaluadas mediante validación cruzada de 3 particiones, totalizando 24 ajustes de modelo. Se utilizó \texttt{n\_jobs=-1} para paralelizar el cómputo sobre todos los núcleos disponibles del entorno de ejecución.
 
\subsection{Validación}
Los hiperparámetros óptimos obtenidos fueron $\texttt{max\_depth}=10$, $\texttt{min\_samples\_split}=2$ y $\texttt{n\_estimators}=100$, alcanzando un F1-Score Macro interno de validación cruzada de $0.8330$.